\lecture{Lecture 4-5: Single pixel ops, img histogram, spatial filtering}

\qa{Define point operations and explain the difference between Logarithmic and Power-Law (Gamma) transformations}{
Point ops are spatial domain image processing techniques where the new value of a pixel depends strictly on the intensity value of that same pixel in the original image, without considering its neighbors. This is expressed as s = T(r), where r is the original intensity, s is the new intensity, and T is the transformation function.

Logarithmic transform $s = c \log(1 + r)$, where $c = \frac{L - 1}{\log(L)}$ expands the values of dark pixels while compressing higher-level values. It is primarily used to compress the dynamic range of an image with huge variations in pixel values, such as displaying the FT spectrum.

Power-law / gamma transform $s = cr^\gamma$, where $c = (L - 1)^{1 - \gamma}$ is more versatile than the log transform as by altering the gamma value, it can both expand dark regions (when $\gamma < 1$) or expand bright regions (when $\gamma > 1$). It is most commonly used for gamma correction to adjust an image's contrast to match the display specs.
}

\qa{Compare and contrast "Contrast Stretching" and "Thresholding".}{
Two techniques that both rely on piecewise-linear transformation functions to manipulate image contrast, but they serve different goals.

Contrast stretching: expands a narrow range of intensity levels in an image to span the full available dynamic range. It is used to correct low-contrast images by making the darks darker and brights brighter, while preserving the relative order of pixel intensities. It is represented as three segments, in which typically the one in middle has a slope coefficient $> 1$ resulting in a stretching of those intensity levels (contrast expansion), while the other two have slope coefficient $< 1$ resulting in an constrast compression.

Thresholding: extreme form of contrast stretching used for image segmentation. It maps all pixel intensities below a specified threshold value to a single dark value (like 0) and all intensities above the threshold to a single bright value (like L-1). Unlike stretching, thresholding irreversibly destroys the original intensity variations, resulting in a binary image.
}

\qa{What is the main limitation of Histogram Equalization, and why might we use Histogram Specification (Matching) instead?}{
The primary limitation of Histogram Equalization is that it is a fully automatic, global operation that aggressively forces the image's histogram to be as flat as possible. While this generally increases contrast, it cannot be controlled and may artificially enhance background noise or wash out details in uniform areas.

Histogram Specification (Matching) addresses this by giving the user control over the output intensity distribution. Unlike Equalization, Specification transforms the input image to match a specific, user-defined target histogram rather than forcing a flat one. This is useful for standardizing images from different sensors, for preventing the over-enhancement caused by standard equalization or for highlighting specific intensity ranges.

The algorithm is a \textbf{global point operation}, meaning the output pixel value depends strictly on the input pixel intensity and the global transformation function (derived from CDFs), not on the neighborhood of pixels. Therefore, it differs from spatial filtering or edge detection, which rely on local operations.

The mathematical procedure consists of three steps:
\begin{enumerate}
    \item \textbf{Equalize Input:} Compute the Cumulative Distribution Function (CDF) of the original image to map original intensities $r$ to intermediate equalized values $s$.
    \item \textbf{Equalize Target:} Compute the CDF of the specified target histogram to map desired output intensities $z$ to the same intermediate equalized values $v$ (where $v \approx s$).
    \item \textbf{Inverse Transform:} Apply the inverse of the target transformation, $z = G^{-1}(s)$, to map the intermediate values $s$ to the final target intensities $z$, resulting in an image with the user-defined histogram shape (neglecting rounding effects).
\end{enumerate}
}

\qa{Explain the concepts of histogram and its equalization. A) Define an image histogram. B) What is the goal of histogram equalization and why is it useful? C) Outline the procedure for deriving the equalization function, providing an example.}{
A) An image histogram is a graphical representation of the distribution of pixel intensities in an image. Specifically, it is a bar graph where the x-axis represents the possible intensity values (0 to L-1) and the y-axis represents the frequency (number of pixels) that assume each specific value. It provides a global description of image's appearance, showing if image is well contrasted, under/overexposed (histogram heavily concentrated on left/right side, respectively). It doesn't contain any spatial information about pixels location in the image.

B) Histogram Equalization is an image enhancement technique used to improve the global contrast of an image. It's useful when an image's pixels intensities are concentrated in a narrow range (e.g. under/overexposed images). The goal is to redistribute the pixel intensities so that they are as evenly spread as possible across the entire available spectrum, resulting in a flatter histogram and higher contrast.

C) That's achieved by computing PDF of original image so to find frequencies for all L intensity levels, then deriving CDF, then map original pixel values $r_k$ to new equalized values. By replacing every pixel with a newly mapped value, we generate an image utilizing the full dynamic range.

\vspace{1em}
Let's assume now to have a 3-bit image ($L=8$), currently very dark (total 10 pixels):

\begin{center}
\begin{tabular}{|c|c|c|c|c|c|}
\hline
\textbf{Original intensity $r_k$} & \textbf{Pixel count $n_k$} & \textbf{PDF ($n_k/10$)} & \textbf{CDF} & \textbf{$(L-1)$CDF} & \textbf{$s_k$} \\
\hline
0 & 2 & 0.2 & 0.2 & 1.4 & 1 \\
1 & 4 & 0.4 & 0.6 & 4.2 & 4 \\
2 & 3 & 0.3 & 0.9 & 6.3 & 6 \\
3 & 1 & 0.1 & 1.0 & 7.0 & 7 \\
4 to 7 & 0 & 0 & 1.0 & 7.0 & 7 \\
\hline
\end{tabular}
\end{center}

Resulting in a brighter image, considering that the new equalized intensities $s_k$ are greater than the original values $r_k$.
}

\qa{Describe the mathematical procedure for performing Histogram Specification.}{
Histogram Specification involves transforming an input image so its intensity distribution matches a specified target PDF. The procedure consists of 3 main steps:
1. Equalize the original image: compute the CDF of the original image to find the transformation $s = T(r)$. This maps the original pixels $r$ to equalized values $s$.
2. Equalize the target histogram: compute CDF of the specified target histogram to find the transformation $v = G(z)$, $z =$ desired output pixel value. This maps the target pixels to equalized values $v$.
3. Apply the inverse transformation: Since both $s, v$ are equalized (flat) distributions, let's assume $s \approx v$ and compute the inverse of the target transformation, $z = G^{-1}(s)$, and map the equalized pixels $s$ from the first step to their final specified values $z$.
}
